% Constructed and checked from the problem; not a geometric proof.
\documentclass[tikz,border=8pt]{standalone}
\begin{document}
\begin{tikzpicture}[scale=1,line cap=round,line join=round]
\coordinate (p0) at (0,0);
\coordinate (p1) at (5,0);
\coordinate (p2) at (0,7);
\coordinate (p3) at (2.5,3.5);
\coordinate (p4) at (2.5,0);
\coordinate (p5) at (2.5,-3.5);
\draw[thick] (p0) -- (p1);
\draw[thick] (p1) -- (p2);
\draw[thick] (p2) -- (p0);
\draw[dashed] (p5) -- (p3);
\draw[dashed] (p0) -- (p5);
\draw[dashed] (p3) -- (p2);
\draw[dashed] (p5) -- (p1);
\draw[dashed] (p1) -- (p3);
\draw[dashed] (p3) -- (p0);
\draw (3.6808327,1.7005948) -- (3.8191673,1.7994052);
\draw (1.1808327,5.2005948) -- (1.3191673,5.2994052);
\draw (1.22025,0.085) -- (1.22025,-0.085);
\draw (1.27975,0.085) -- (1.27975,-0.085);
\draw (3.72025,0.085) -- (3.72025,-0.085);
\draw (3.77975,0.085) -- (3.77975,-0.085);
\fill (p0) circle (1.1pt) node[fill=white,inner sep=1pt,below left] {$A$};
\fill (p1) circle (1.1pt) node[fill=white,inner sep=1pt,below right] {$B$};
\fill (p2) circle (1.1pt) node[fill=white,inner sep=1pt,above left] {$C$};
\fill (p3) circle (1.1pt) node[fill=white,inner sep=1pt,above right] {$M$};
\fill (p4) circle (1.1pt) node[fill=white,inner sep=1pt,below right] {$D$};
\fill (p5) circle (1.1pt) node[fill=white,inner sep=1pt,below right] {$E$};
\end{tikzpicture}
\end{document}
